\section{Literature and Position}
\label{sec:lit}

This section positions the paper in four literatures --- disclosure-threshold
theory, the theory of activism and the market for corporate control, the
structural quantification of activism, and the measurement of takeover gains ---
and then states the reform and the novelty claim. The characterizations below
are disciplined by a machine-checked reading of the full texts; where a claim
rests on a verbatim quotation, the quotation is logged in the accompanying
verification report. The through-line is that the activist-as-intermediary
position is occupied at the theory level, while two margins are open: the filing
\emph{window} as a time margin of stake formation, and \emph{realized} evidence
on the February 2024 acceleration.

\subsection{Disclosure-threshold theory}
\label{sec:lit-threshold}

\textcite{grossmanhart1980disclosure} is the classical antecedent: mandatory
disclosure of a building position raises the acquirer's cost of stealth and so
trades off the protection of target shareholders against the disciplining of
management. The modern theoretical treatment is
\textcite{ordonezcalafibernhardt2022}, who study a blockholder disclosure
threshold as a \emph{static stake-size cap}: in a one-shot Kyle-style dealership
trade, a competitive market maker posts prices against net order flow, and the
activist chooses a single position once. Their policy instrument is an ownership
level --- a shareholder must disclose once holdings reach a given fraction of
voting rights --- not a time margin. There is no accumulation period, no filing
deadline, and no second agent representing a potential bidder; the welfare result
is the familiar transparency-versus-discipline trade-off, in which uninformed
investors favor lower thresholds when their trading losses against the informed
activist are small relative to the activist's disciplining benefit. Critically
for us, their Section~V explicitly gestures at objective-specific thresholds ---
that the right disclosure rule may depend on what the activist is trying to do,
including sale-seeking campaigns --- but leaves this as an open question. Answering
it, with the filing window rather than the stake level as the lever, is part of
what this paper does.

\textcite{bonettiduroormazabal2020} is the structurally closest \emph{design}
paper, and it is empirical rather than theoretical. Using the staggered
country-level implementation of the E.U.\ Transparency Directive (2007--2009) as a
difference-in-differences shock, they show that tighter ownership disclosure
\emph{shrinks toeholds, reduces the number of control acquisitions, and moves
announcement returns in opposite directions for target and acquirer}. Their
setting is ordinary corporate bidders building pre-bid toeholds, not hedge-fund
activist 13D campaigns, and their reform bundles a lower ownership threshold with a
shorter notification deadline, so it is not an isolated window shock the way the
2024 Schedule~13D change is. The reason to engage them head-on is the
\emph{sign} of the bidder-side channel. Building on \textcite{fishman1988}, their
bidding-competition mechanism has disclosure \emph{deterring} rival bidders: once
alerted, competitors can make competing bids that raise the acquirer's cost. That
is the opposite sign from the certification channel this paper emphasizes, in
which earlier disclosure \emph{attracts} bidders by signaling that an
activist-flagged target is available for sale. The two signs are not
contradictory; they are the two sides of a single elasticity. Our
certification-versus-deterrence result (Proposition~R4) nests both: the bid
hazard rises after earlier disclosure if and only if the certification elasticity
exceeds the deterrence elasticity, and the observable heterogeneity that signs it
(sale-intent filings, opaque assets with few natural buyers, high industry
merger intensity) is precisely what separates an activist-flagged sale from the
ordinary toehold acquisitions \textcite{bonettiduroormazabal2020} study.

\subsection{Activism and the market for corporate control}
\label{sec:lit-activism-theory}

The theoretical position this paper builds on --- the activist as an intermediary
between an underperforming target and the takeover market --- is already occupied,
and we say so plainly. \textcite{burkartlee2022} make it their headline: in their
words, ``activism is most efficient when it brokers, rather than substitutes for,
takeovers.'' Their activist runs a proxy-style campaign for control of the board
and then uses the board's merger authority to sell the whole firm to a
third-party bidder, mitigating both the Grossman--Hart and Jensen--Meckling
free-rider problems. But the toehold is exogenous and the model is entirely
post-disclosure: there is no stake accumulation and no window. Their concluding
remarks explicitly defer ``the acquisition of the toehold in anonymous,
predisclosure markets'' to future work --- the margin we open.

\textcite{corumlevit2019} occupy the bidder-discovery side of the same position.
In their model bidders ``interpret the presence of an activist as a signal that
the target is available for sale,'' so the activist's presence puts the target in
play and raises the probability of a takeover. It is important to be precise
about what they do and do not model, because an earlier draft overstated this.
They \emph{do} model secondary-market inference: the activist builds a position
by trading against a \textcite{kyle1985} market maker who learns from order flow,
and the activist's ownership then becomes public through a Schedule~13D
publication step. What they do \emph{not} model is a disclosure \emph{regime} ---
the publication step is instantaneous, with no threshold design and no filing
window or deadline to vary. Our contribution relative to them is therefore narrow
and honest: the window-as-deadline time margin, and the realized reform, not the
claim that they ignore trading. Their footnote~29 also already delivers the sign
often credited to later structural work: when the activist's proxy threat is
credible, a takeover is more likely but the bidder pays a lower conditional
premium, because the activist weakens the board's negotiating position.

\textcite{cetemen2026} is a continuous-time leader--follower model of how a
leader activist's trades tacitly coordinate followers, producing wolf-pack
accumulation without formal agreement. It is pure theory about coordination
through trading, and it is \emph{not} silent on the reform --- a footnote cites
the February 2024 acceleration (ten calendar days to five business days) as
institutional background --- but it does not incorporate the window as a variable
or analyze it. The honest claim is therefore that no theory paper gives the window
\emph{analytical} treatment, not that the reform goes unmentioned.

Two strategic-trading papers are the benchmarks our reduced-form accumulation cost
shortcuts. \textcite{back2018} endogenize the activist's optimal trading and
liquidity, and \textcite{collindufresnefos2016} solve the problem of an informed
trader with long-lived private information revealed at a known date, trading
against stochastic liquidity. A static quadratic cost
$C(x;d,L)=(\lambda_L(L)/2)\,x^{2}/d$ is a fair reduced-form summary of the
\emph{level} effects these models deliver --- more liquidity and a longer horizon
lower the marginal cost of a given accumulation --- but it collapses away their
two dynamic contributions: endogenous intra-window timing (the insider trades more
aggressively when noise trading is high) and stochastic, path-dependent
liquidity. We use the reduced form to generate the exposure measure and
organizing predictions, and we are explicit that a full dynamic treatment would
add a timing margin the static cost cannot represent. The empirical reality that
motivates a reduced form is \textcite{collindufresnefos2015}: on an average day
when Schedule~13D filers trade, measured price impact is almost 30\% lower than
the sample average, so standard liquidity proxies do not cleanly detect informed
window trading --- a caution for any design that uses pre-filing liquidity as a
control.

The older blockholder literature supplies the background on exit, voice, and
liquidity. \textcite{maug1998} shows that liquidity raises the value of the
exit option and so trades off against monitoring; \textcite{kahnwinton1998}
argue that liquid markets can substitute speculation for intervention;
\textcite{edmansfangzur2013} document that liquidity shifts blockholders away from
voice and toward governance through exit; and \textcite{edmansholderness2017}
survey the theory and evidence on blockholders. None speaks to a disclosure
window or to takeover-gain timing.

\subsection{Quantifying activism: structural evidence}
\label{sec:lit-structural}

The structural literature quantifies campaigns but does not exploit disclosure
timing. \textcite{johnsonswem2021} estimate a dynamic reputation model in which
an activist's track record for proxy fighting shapes targets' responses across a
sequence of campaigns. Accumulation is collapsed into a single reduced-form
lognormal campaign-cost draw --- round-trip liquidity costs are named as a verbal
component of that one scalar but never separately estimated --- and the filing
window appears exactly once, in an institutional footnote describing the SEC's
ten-day rule. There is no mapping of the window to a noise or variance parameter
and no counterfactual that varies filing delay; the model in fact has no
separable disclosure-timing choice, because filing the 13D is definitionally
simultaneous with campaign initiation.

\textcite{albuquerquefosschroth2022} study the choice between Schedule~13D and
Schedule~13G and decompose the average 13D announcement return of 6.34\% into
treatment (75.2\%), stock-picking (12.2\%), and sample-selection (12.6\%)
components --- a capitalization-at-disclosure result at the filing leg. Their
activism cost is a single scalar held constant across filings, and they offer no
takeover-premium counterfactual; an earlier project document misattributed the
premium result below to them, and we correct that attribution here.

\textcite{celentanolevine2025} estimate a simulated-method-of-moments model on
2001--2019 data and report that, in takeovers with an activist present, ``the bid
premium is 13.7 percent (5.2 percentage points) lower than if the activist had
not been present,'' with the mechanism that the activist ``lowers the negotiating
position and price demanded by the board in negotiations'' --- reduced board
bargaining leverage, not a liquidity or disclosure channel. Crucially, they
themselves assert the capitalization reading of their own number: ``the
pre-announcement share price used to calculate the bid premium incorporates the
activist's signal about a potential acquirer.'' Because the premium's denominator
already capitalizes the campaign, a lower measured premium need not mean less
value to target shareholders. This is exactly why we treat capitalization as a
measurement lemma rather than a headline: \textcite{celentanolevine2025} own the
interpretation, and their model treats the activist's decision and information as
instantaneously common knowledge, with no disclosure timing anywhere, leaving the
window and the realized reform open.

\textcite{gantchev2013} estimates that the average activist campaign costs about
\$10.5 million, roughly half of it in the proxy stage. These are the costs of
contested campaign \emph{tactics} --- proxy-fight, legal, and advisory
expenditure --- and are conceptually distinct from the accumulation cost this
paper's model represents, the market-impact cost of quietly building a
pre-disclosure stake. We cite the figure for scale, not as a calibration of our
cost object.

\subsection{Takeovers, run-ups, and the measurement of gains}
\label{sec:lit-measurement}

The gain anatomy borrows the standard run-up/markup decomposition.
\textcite{schwert1996} writes the premium as the sum of a pre-bid run-up,
measured as the target's cumulative abnormal return from day $-42$ to day $-1$,
and a post-announcement markup. His headline finding, however, runs
\emph{against} substitution: he reports ``little substitution between the runup
and the markup,'' with at least two-thirds of any run-up \emph{added} to the total
price the bidder pays. The substitution evidence our measurement lemma actually
needs comes from \textcite{bett2014}: once the run-up is adjusted for the target's
own stand-alone value change, isolating the deal-synergy component, the slope
moves from $-0.09$ (following Schwert's unadjusted method) to $-0.39$, indicating
``a substantially greater degree of substitution'' and rejecting the mechanical
costly-feedback loop. The run-up/markup window convention is standardized in the
survey of \textcite{eckbomalenkothorburn2025}, which uses a run-up window of day
$-41$ to day $-2$ and a separate three-day announcement window. Our three-part
split --- filing, run-up, and markup --- adds an activist-filing-triggered leg
that has no counterpart in the two-part bid-anchored convention, and we justify it
on its own terms.

The activism--takeover link is a stylized fact. \textcite{greenwoodschor2009}
show that the large abnormal returns to 13D activism are ``largely explained by
the ability of activists to force target firms into a takeover'': announcement
and long-run returns are high only for targets ultimately acquired and
indistinguishable from zero otherwise, and activism raises takeover probability by
about eleven percentage points. \textcite{brav2008} report that selling the
company is a stated objective in roughly 14--18\% of campaigns and that about
15\% of activist targets are ultimately acquired. \textcite{boysongantchevshivdasani2017}
study activism mergers directly and find that activist targets are far more likely
to receive bids and to merge; our characterization rests on an author-posted
summary rather than the primary text, so we keep it to this one conservative
sentence and do not build our sale-intent codebook on it. None of these papers
touches disclosure timing.

\subsection{The reform and the novelty claim}
\label{sec:lit-reform}

The institutional shock is the SEC's modernization of beneficial-ownership
reporting \parencite{sec2023modernization}. The final rule shortens the initial
Schedule~13D filing deadline from ten calendar days to five business days,
effective February~5, 2024; requires amendments within two business days of a
material change; and mandates XML-structured filing from December~18, 2024. Two
features matter for the design. First, the rule \emph{dropped} proposed Rule
13d-3(e): the Commission declined to deem holders of cash-settled derivative
securities beneficial owners of the reference class, issuing guidance and an
Item~6 disclosure clarification instead. That decision leaves a substitution hole
--- activists can migrate economic exposure into cash-settled derivatives without
triggering a 13D --- which is the second half of our chilling/substitution
headline. Second, the SEC's own baseline shows the reform binds a minority of
campaigns: roughly 29\% of all 2022 initial filings (about 41\% of those filed on
time) were already made within five business days, about 80\% of filers had completed accumulation by the new
deadline, and 97\% of campaigns had reached 90\% of their reported stake by it.
The staff also projected foregone shareholder value of about \$810 million per
year under the five-day rule --- a projection this paper adjudicates with realized
data. The February 2025 staff guidance \parencite{seccdi2025} clarifies when such
engagement makes the shorter Schedule~13G unavailable (calls for a sale or
restructuring, alternative director nominees, or conditional voting support),
which bears on how sale-intent campaigns are routed between the two schedules.

The only prior study of the acceleration itself is \textcite{polk2024}, a
pre-reform projection on 9,685 initial 13D filings from 2001--2022 that compares
investors who voluntarily filed within five days against those who waited the full
ten. Its outcomes are abnormal returns and trading volume around the trigger and
filing dates only: it contains no takeover, premium, bidder-entry, or
campaign-resolution variables, and its sample ends more than a year before the
reform took effect. A systematic novelty scan run on 2026-08-02 --- spanning SSRN,
NBER, arXiv, the SEC/DERA record, forthcoming top-journal pages, and author
research pages --- found \emph{no} study of realized post-February-2024 Schedule
13D outcomes. The closest adjacent work, \textcite{celentanolevine2025}, is not
built around the reform and uses no post-reform data. The realized evaluation is
genuinely open.

\subsection{Position}
\label{sec:lit-position}

The activist-as-intermediary position is taken at the theory level, and the
capitalization reading of a lower conditional premium is asserted by its authors;
we do not reclaim either as new. What remains open are two margins. The first is
the filing window as a \emph{time margin} of stake formation: disclosure-threshold
theory fixes a static stake-size cap, the brokerage and solicitation models take
the toehold as given or publish it instantaneously, and the structural literature
collapses accumulation into a scalar cost draw, so no paper makes the legal
deadline $d$ a margin that binds only for a model-generated subset of targets.
The second is \emph{realized} evidence on the February 2024 acceleration: the
existing study is a pre-reform projection with returns and volume only. This
paper supplies both --- a deadline-accumulation model used as a lens to generate
the predetermined required-trading-days exposure
$\mathrm{RTD}_i = (0.05 \times S_{\mathrm{out},i})/(\phi \times \mathrm{ADV}_i)$
and to organize predictions, and a dose-response design on that exposure as the
evidence. The event-time-invariant anatomy of target gains (filing, run-up,
markup) is the measurement contribution that disciplines the capitalization
lemma; the model is the lens, not the paper's claim. Table~\ref{tab:position}
records where each paper sits.

\begin{table}[t]
\centering
\small
\caption{Positioning. Columns: (Window) the filing window modeled as a time
margin of stake formation; (Realized) realized evidence on the February 2024
Schedule~13D acceleration; (Anatomy) a decomposition of target gains into
filing/run-up/markup; (13D) an activist Schedule~13D setting.}
\label{tab:position}
\begin{tabular}{@{}lcccc@{}}
\toprule
 & Window & Realized & Anatomy & 13D \\
\midrule\relax
\textcite{ordonezcalafibernhardt2022} & no  & no  & no      & yes \\\relax
\textcite{burkartlee2022}             & no  & no  & no      & yes \\\relax
\textcite{corumlevit2019}             & no  & no  & no    & yes \\\relax
\textcite{johnsonswem2021}            & no  & no  & no      & yes \\\relax
\textcite{albuquerquefosschroth2022}  & no  & no  & partial & yes \\\relax
\textcite{celentanolevine2025}        & no  & no  & partial & yes \\\relax
\textcite{bonettiduroormazabal2020}   & partial & no & no    & no \\\relax
\textcite{polk2024}                   & partial & no & no    & yes \\\relax
\textcite{cetemen2026}                & no  & no  & no      & yes \\
\midrule
\emph{This paper}                     & yes & yes & yes     & yes \\
\bottomrule
\end{tabular}
\end{table}
